\input{preamble}

\begin{document}

\title{Geometric prequantization}
\maketitle

\phantomsection
\label{section-phantom}

\tableofcontents

\section{Overview}
\label{section-overview}

\noindent
These notes are about the first step in
geometric quantization.  The starting point
is
a symplectic manifold $(M,\omega)$,
regarded as
the phase space of a classical mechanical
system.  The goal is to attach to it a
Hilbert
space and a rule sending classical
observables
$f \in C^\infty(M)$ to operators on that
Hilbert space.

\medskip\noindent
The construction treated here is only
prequantization.  It gives the correct formal
commutator rule, but the resulting Hilbert
space is usually too large.  The later step,
not treated here, is to impose a
polarization.

\medskip\noindent
We use the following convention.  The
Hamiltonian vector field of $f$ is defined by
$$
\iota_{X_f}\omega=-df,
$$
and the Poisson bracket is
$$
\{f,g\}=-\omega(X_f,X_g).
$$
With this convention $X_f(g)=-\{f,g\}$.

\section{The Dirac rule}
\label{section-dirac-rule}

\noindent
In classical mechanics an observable is a
smooth function on the phase space.  The
Hamiltonian flow of $H$ evolves another
observable $f$ by
$$
\frac{d}{dt}f=\{f,H\}.
$$
Thus the Poisson bracket is the classical
operation governing the time evolution of
observables.

\medskip\noindent
In quantum mechanics observables are
self-adjoint operators on a Hilbert space
$\mathcal H$.  If $H_q$ is the quantum
Hamiltonian, the Schrodinger equation is
$$
i\hbar\frac{d\psi}{dt}=H_q\psi.
$$
For a time-independent observable $A$, the
expected value in the state $\psi$ is
$$
\left\langle{} A \right\rangle{}_\psi
=\left\langle{} A\psi,\psi \right\rangle{}.
$$
The commutator controls its evolution:
$$
\frac{d}{dt}
\left\langle{} A \right\rangle{}_\psi
=
\frac{1}{i\hbar}
\left\langle{} [A,H_q]\psi,\psi
\right\rangle{}.
$$
This is the basic reason the commutator is
the quantum analogue of the Poisson bracket.

\begin{definition}
\label{definition-dirac-prequantization}
Let $(M,\omega)$ be a symplectic manifold.
A {\it prequantization} of $M$ consists of a
Hilbert space $\mathcal H$ and a linear map
$$
Q:C^\infty(M,\mathbb R)\longrightarrow
\operatorname{End}(\mathcal H)
$$
such that
\begin{enumerate}
\item $Q(1)=\operatorname{Id}$, and
\item
$$
Q(\{f,g\})=\frac{1}{i\hbar}[Q(f),Q(g)]
$$
for all $f,g\in C^\infty(M,\mathbb R)$.
\end{enumerate}
\end{definition}

\noindent
A full quantization should also include a
minimality condition saying that a complete
set of commuting classical observables goes
to a complete set of commuting quantum
observables.  This is where
prequantization is
not yet enough.

\section{Hermitian line bundles}
\label{section-hermitian-line-bundles}

\begin{definition}
\label{definition-hermitian-line-bundle}
A {\it Hermitian line bundle} over $M$ is a
complex line bundle $L\to M$ together with a
smooth Hermitian product $h$ on the fibers.
Thus $h(s,t)$ is a smooth complex-valued
function whenever $s,t$ are smooth sections.
\end{definition}

\noindent
We write
$$
\Omega^k(M,L)=
\Gamma(M,\Lambda^kT^*M\otimes L)
$$
for the space of $L$-valued differential
forms.

\begin{definition}
\label{definition-connection-line-bundle}
A {\it connection} on $L$ is a linear map
$$
\nabla:\Gamma(M,L)\longrightarrow
\Omega^1(M,L)
$$
such that
$$
\nabla(fs)=df\otimes s+f\nabla s
$$
for all $f\in C^\infty(M)$ and all sections
$s$ of $L$.
For a vector field $X$ we write
$$
\nabla_Xs=\iota_X\nabla s.
$$
\end{definition}

\begin{definition}
\label{definition-unitary-connection}
A connection $\nabla$ on a Hermitian line
bundle $(L,h)$ is {\it unitary} if
$$
dh(s,t)=h(\nabla s,t)+h(s,\nabla t)
$$
for all smooth sections $s,t$.
\end{definition}

\noindent
A connection extends to $L$-valued forms by
$$
\nabla(\alpha\otimes s)
=d\alpha\otimes s-\alpha\wedge \nabla s.
$$
Then $\nabla^2$ is $C^\infty(M)$-linear, so
there is a complex 2-form $\Omega_\nabla$
such that
$$
\nabla^2s=\Omega_\nabla s.
$$
This form is the curvature of $\nabla$.

\begin{lemma}
\label{lemma-local-curvature-line-bundle}
Let $e$ be a nowhere vanishing local section
of $L$ and suppose
$$
\nabla e=\theta\otimes e.
$$
Then, on the same open set,
$$
\Omega_\nabla=d\theta.
$$
If $e$ is unitary and $\nabla$ is unitary,
then $\theta$ and $\Omega_\nabla$ are purely
imaginary forms.
\end{lemma}

\begin{proof}
For vector fields $X,Y$ we have
$$
\begin{aligned}
\nabla_X\nabla_Ye
&=\nabla_X(\theta(Y)e)\\
&=X(\theta(Y))e+	heta(Y)\theta(X)e.
\end{aligned}
$$
Subtracting the same expression with $X$ and
$Y$ interchanged, and subtracting
$\nabla_{[X,Y]}e$, gives
$$
\Omega_\nabla(X,Y)e
=d\theta(X,Y)e.
$$
Hence $\Omega_\nabla=d\theta$.

If $e$ is unitary, then $h(e,e)=1$.  Since
$\nabla$ is unitary,
$$
0=dh(e,e)=h(\nabla e,e)+h(e,\nabla e)
=\theta+\overline\theta.
$$
Thus $\theta$ is purely imaginary, and so is
$d\theta$.
\end{proof}

\begin{definition}
\label{definition-first-chern-class}
The {\it first Chern class} of a Hermitian
line bundle with unitary connection is
$$
c_1(L)=\left[\frac{1}{2\pi i}
\Omega_\nabla\right]\in H^2(M,\mathbb R).
$$
\end{definition}

\noindent
Chern--Weil theory says this class is
independent of the chosen unitary connection.
Moreover it is integral, i.e. it lies in the
image of $H^2(M,\mathbb Z)$ in real
cohomology.

\begin{theorem}[Weil integrality theorem]
\label{theorem-weil-integrality}
\begin{reference}
See \cite[Lecture 12]{wang} and
\cite[Section 8.3]{wood}.
\end{reference}
Let $\beta$ be a closed real 2-form on $M$
whose de Rham cohomology class is integral.
Then there exists a Hermitian line bundle
$L\to M$ with unitary connection $\nabla$
such that
$$
\Omega_\nabla=2\pi i\beta.
$$
Equivalently, $c_1(L)=[\beta]$.
\end{theorem}

\begin{proof}
We omit the proof.  It is the standard
Cech--de Rham construction of a line bundle
from an integral cocycle, followed by the
construction of a compatible unitary
connection.
\end{proof}

\section{Prequantum line bundles}
\label{section-prequantum-line-bundles}

\begin{definition}
\label{definition-prequantizable}
A symplectic manifold $(M,\omega)$ is
{\it prequantizable} if
$$
\left[\frac{\omega}{2\pi\hbar}\right]
\in H^2(M,\mathbb Z).
$$
A {\it prequantum line bundle} is a
Hermitian line bundle $(L,h)$ with a unitary
connection $\nabla$ whose curvature is
$$
\Omega_\nabla=\frac{\omega}{i\hbar}.
$$
\end{definition}

\noindent
The sign in the integrality condition is
irrelevant: if a class is integral, so is its
negative.  By Theorem
\ref{theorem-weil-integrality}, the
integrality condition is exactly what is
needed for the existence of a prequantum line
bundle.

\medskip\noindent
When $M$ is compact, the prequantum Hilbert
space is the $L^2$-completion of the smooth
sections of $L$ with respect to
$$
\left\langle{} s,t\right\rangle{}
=
\frac{1}{(2\pi\hbar)^n}
\int_M h(s,t)\frac{\omega^n}{n!}.
$$
The measure $\omega^n/n!$ is the Liouville
measure.

\begin{definition}
\label{definition-prequantum-operator}
Let $f\in C^\infty(M,\mathbb R)$.
The {\it prequantum operator} associated to
$f$ is
$$
Q(f)=-i\hbar\nabla_{X_f}+m_f,
$$
where $m_f$ denotes multiplication by $f$.
\end{definition}

\section{The Kostant--Souriau formula}
\label{section-kostant-souriau-formula}

\begin{lemma}
\label{lemma-hamiltonian-volume-preserving}
For every $f\in C^\infty(M)$,
$$
\mathcal L_{X_f}\omega=0
$$
and hence
$$
\mathcal L_{X_f}\omega^n=0.
$$
\end{lemma}

\begin{proof}
By Cartan's formula and the definition of
$X_f$,
$$
\mathcal L_{X_f}\omega
=d\iota_{X_f}\omega+
\iota_{X_f}d\omega
=-d^2f+0=0.
$$
The statement for $\omega^n$ follows by the
Leibniz rule for the Lie derivative.
\end{proof}

\begin{lemma}
\label{lemma-integral-poisson-bracket-zero}
Assume that $M$ is compact without boundary.
Then, for all $f,g\in C^\infty(M)$,
$$
\int_M \{f,g\}\frac{\omega^n}{n!}=0.
$$
\end{lemma}

\begin{proof}
Since $X_f(g)=-\{f,g\}$ and
$\mathcal L_{X_f}\omega^n=0$, we have
$$
\{f,g\}\omega^n
=-\mathcal L_{X_f}(g\omega^n).
$$
By Cartan's formula,
$$
\mathcal L_{X_f}(g\omega^n)
=d\iota_{X_f}(g\omega^n)+
\iota_{X_f}d(g\omega^n).
$$
The second term vanishes for degree reasons,
because $g\omega^n$ is a top-degree form.
Thus $\{f,g\}\omega^n$ is exact, and the
integral vanishes by Stokes' theorem.
\end{proof}

\begin{proposition}
\label{proposition-prequantum-self-adjoint}
Assume that $M$ is compact without boundary.
For real $f$, the operator $Q(f)$ is formally
self-adjoint on smooth sections.
\end{proposition}

\begin{proof}
The multiplication operator $m_f$ is formally
self-adjoint because $f$ is real.  It remains
to check $-i\hbar\nabla_{X_f}$.
By compatibility of $\nabla$ with $h$,
$$
X_fh(s,t)=h(\nabla_{X_f}s,t)
+h(s,\nabla_{X_f}t).
$$
Integrating and using Lemma
\ref{lemma-integral-poisson-bracket-zero} for
real and imaginary parts of $h(s,t)$, we get
$$
\int_M X_fh(s,t)\frac{\omega^n}{n!}=0.
$$
Hence
$$
\int_M h(\nabla_{X_f}s,t)\frac{\omega^n}{n!}
=-
\int_M h(s,\nabla_{X_f}t)\frac{\omega^n}{n!}.
$$
Thus $\nabla_{X_f}$ is formally skew-adjoint,
and $-i\hbar\nabla_{X_f}$ is formally
self-adjoint.
\end{proof}

\begin{theorem}[Kostant--Souriau]
\label{theorem-kostant-souriau}
\begin{reference}
See \cite[Lecture 12]{wang} and
\cite[Chapter 23]{hallq}.
\end{reference}
Let $(M,\omega)$ be prequantizable, and let
$(L,h,\nabla)$ be a prequantum line bundle.
Then
$$
Q(\{f,g\})=\frac{1}{i\hbar}[Q(f),Q(g)]
$$
for all $f,g\in C^\infty(M,\mathbb R)$.
Thus $f\mapsto Q(f)$ is a prequantization in
the sense of Definition
\ref{definition-dirac-prequantization}.
\end{theorem}

\begin{proof}
Let $h=\{f,g\}$.  First,
$$
{}[X_f,X_g]=-X_h.
$$
Indeed,
$$
\iota_{[X_f,X_g]}\omega
=\mathcal L_{X_f}\iota_{X_g}\omega
=-\mathcal L_{X_f}dg
=-d(X_fg)=dh,
$$
because $X_f(g)=-h$.  Since
$\iota_{X_h}\omega=-dh$, this gives
$[X_f,X_g]=-X_h$.

Now compute the commutator on smooth
sections.  We use
$$
{}[\nabla_X,m_g]=m_{X(g)}
$$
and the curvature identity
$$
{}[\nabla_X,\nabla_Y]
=\nabla_{[X,Y]}+\Omega_\nabla(X,Y).
$$
Therefore
$$
\begin{aligned}
{}[Q(f),Q(g)]
&=(-i\hbar)^2[\nabla_{X_f},\nabla_{X_g}]\\
&\quad -i\hbar m_{X_f(g)}
+i\hbar m_{X_g(f)}.
\end{aligned}
$$
Since $X_f(g)=-h$ and $X_g(f)=h$, the last
two terms give $2i\hbar m_h$.  For the
first term we use the prequantum curvature:
$$
\begin{aligned}
{}[\nabla_{X_f},\nabla_{X_g}]
&=\nabla_{-X_h}
+\frac{1}{i\hbar}\omega(X_f,X_g)\\
&=-\nabla_{X_h}-\frac{1}{i\hbar}h.
\end{aligned}
$$
Thus
$$
\begin{aligned}
{}[Q(f),Q(g)]
&=(-i\hbar)^2
\left(-\nabla_{X_h}
-\frac{1}{i\hbar}m_h\right)
+2i\hbar m_h\\
&=\hbar^2\nabla_{X_h}+i\hbar m_h\\
&=i\hbar(-i\hbar\nabla_{X_h}+m_h)\\
&=i\hbar Q(h).
\end{aligned}
$$
Dividing by $i\hbar$ proves the formula.
Linearity is immediate from the definition of
$Q$, and $Q(1)=\operatorname{Id}$ because
$X_1=0$.
\end{proof}

\section{What remains}
\label{section-what-remains}

\noindent
The prequantum Hilbert space is usually too
large.  Geometric quantization reduces it by
choosing a polarization and keeping only
sections that are covariantly constant along
that polarization.  For example, on Kahler
manifolds this leads to holomorphic sections
of the prequantum line bundle.

\medskip\noindent
There is also a structural obstruction to
asking for a perfect quantization of all
classical observables: the Groenewold--van
Hove theorem says that the full Dirac rule
cannot hold for all polynomial observables on
standard phase space.  This is one reason why
prequantization is only the first step.

\bibliography{my}
\bibliographystyle{amsalpha}

\end{document}
